\documentclass[a4paper,12pt,twoside, openright]{memoir}
\usepackage[left=4cm,right=3cm]{geometry}

% Castellano
\usepackage[spanish,es-tabla]{babel}
\selectlanguage{spanish}
\usepackage[utf8]{inputenc}
\usepackage[T1]{fontenc}
\usepackage{lmodern} % scalable font
\usepackage{microtype}
\usepackage{placeins}

\RequirePackage{booktabs}
\RequirePackage[table]{xcolor}
\RequirePackage{xtab}
\RequirePackage{multirow}

% Links
\PassOptionsToPackage{hyphens}{url}\usepackage[colorlinks]{hyperref}
\hypersetup{
	allcolors = {red}
}

% Ecuaciones
\usepackage{amsmath}

% Rutas de fichero / paquete
\newcommand{\ruta}[1]{{\sffamily #1}}

% Párrafos
\nonzeroparskip

% Huérfanas y viudas
\widowpenalty100000
\clubpenalty100000

% Evitar solapes en el header
\nouppercaseheads

% Imagenes
\usepackage{graphicx}
\newcommand{\imagen}[2]{
	\begin{figure}[!ht]
		\centering
		\includegraphics[width=0.9\textwidth]{#1}
		\caption{#2}\label{fig:#1}
	\end{figure}
	\FloatBarrier
}

\usepackage{graphicx}
\newcommand{\imagennormal}[2]{
	\begin{figure}[!ht]
		\centering
		\includegraphics[width=\textwidth]{#1}
		\caption{#2}\label{fig:#1}
	\end{figure}
	\FloatBarrier
}

\newcommand{\imagenflotante}[2]{
	\begin{figure}%[!h]
		\centering
		\includegraphics[width=0.9\textwidth]{#1}
		\caption{#2}\label{fig:#1}
	\end{figure}
}



% El comando \figura nos permite insertar figuras comodamente, y utilizando
% siempre el mismo formato. Los parametros son:
% 1 -> Porcentaje del ancho de página que ocupará la figura (de 0 a 1)
% 2 --> Fichero de la imagen
% 3 --> Texto a pie de imagen
% 4 --> Etiqueta (label) para referencias
% 5 --> Opciones que queramos pasarle al \includegraphics
% 6 --> Opciones de posicionamiento a pasarle a \begin{figure}
\newcommand{\figuraConPosicion}[6]{%
  \setlength{\anchoFloat}{#1\textwidth}%
  \addtolength{\anchoFloat}{-4\fboxsep}%
  \setlength{\anchoFigura}{\anchoFloat}%
  \begin{figure}[#6]
    \begin{center}%
      \Ovalbox{%
        \begin{minipage}{\anchoFloat}%
          \begin{center}%
            \includegraphics[width=\anchoFigura,#5]{#2}%
            \caption{#3}%
            \label{#4}%
          \end{center}%
        \end{minipage}
      }%
    \end{center}%
  \end{figure}%
}

%
% Comando para incluir imágenes en formato apaisado (sin marco).
\newcommand{\figuraApaisadaSinMarco}[5]{%
  \begin{figure}%
    \begin{center}%
    \includegraphics[angle=90,height=#1\textheight,#5]{#2}%
    \caption{#3}%
    \label{#4}%
    \end{center}%
  \end{figure}%
}
% Para las tablas
\newcommand{\otoprule}{\midrule [\heavyrulewidth]}
%
% Nuevo comando para tablas pequeñas (menos de una página).
\newcommand{\tablaSmall}[5]{%
 \begin{table}
  \begin{center}
   \rowcolors {2}{gray!35}{}
   \begin{tabular}{#2}
    \toprule
    #4
    \otoprule
    #5
    \bottomrule
   \end{tabular}
   \caption{#1}
   \label{tabla:#3}
  \end{center}
 \end{table}
}

%
%Para el float H de tablaSmallSinColores
\usepackage{float}

%
% Nuevo comando para tablas pequeñas (menos de una página).
\newcommand{\tablaSmallSinColores}[5]{%
 \begin{table}[H]
  \begin{center}
   \begin{tabular}{#2}
    \toprule
    #4
    \otoprule
    #5
    \bottomrule
   \end{tabular}
   \caption{#1}
   \label{tabla:#3}
  \end{center}
 \end{table}
}

\newcommand{\tablaApaisadaSmall}[5]{%
\begin{landscape}
  \begin{table}
   \begin{center}
    \rowcolors {2}{gray!35}{}
    \begin{tabular}{#2}
     \toprule
     #4
     \otoprule
     #5
     \bottomrule
    \end{tabular}
    \caption{#1}
    \label{tabla:#3}
   \end{center}
  \end{table}
\end{landscape}
}

%
% Nuevo comando para tablas grandes con cabecera y filas alternas coloreadas en gris.
\newcommand{\tabla}[6]{%
  \begin{center}
    \tablefirsthead{
      \toprule
      #5
      \otoprule
    }
    \tablehead{
      \multicolumn{#3}{l}{\small\sl continúa desde la página anterior}\\
      \toprule
      #5
      \otoprule
    }
    \tabletail{
      \hline
      \multicolumn{#3}{r}{\small\sl continúa en la página siguiente}\\
    }
    \tablelasttail{
      \hline
    }
    \bottomcaption{#1}
    \rowcolors {2}{gray!35}{}
    \begin{xtabular}{#2}
      #6
      \bottomrule
    \end{xtabular}
    \label{tabla:#4}
  \end{center}
}

%
% Nuevo comando para tablas grandes con cabecera.
\newcommand{\tablaSinColores}[6]{%
  \begin{center}
    \tablefirsthead{
      \toprule
      #5
      \otoprule
    }
    \tablehead{
      \multicolumn{#3}{l}{\small\sl continúa desde la página anterior}\\
      \toprule
      #5
      \otoprule
    }
    \tabletail{
      \hline
      \multicolumn{#3}{r}{\small\sl continúa en la página siguiente}\\
    }
    \tablelasttail{
      \hline
    }
    \bottomcaption{#1}
    \begin{xtabular}{#2}
      #6
      \bottomrule
    \end{xtabular}
    \label{tabla:#4}
  \end{center}
}

%
% Nuevo comando para tablas grandes sin cabecera.
\newcommand{\tablaSinCabecera}[5]{%
  \begin{center}
    \tablefirsthead{
      \toprule
    }
    \tablehead{
      \multicolumn{#3}{l}{\small\sl continúa desde la página anterior}\\
      \hline
    }
    \tabletail{
      \hline
      \multicolumn{#3}{r}{\small\sl continúa en la página siguiente}\\
    }
    \tablelasttail{
      \hline
    }
    \bottomcaption{#1}
  \begin{xtabular}{#2}
    #5
   \bottomrule
  \end{xtabular}
  \label{tabla:#4}
  \end{center}
}



\definecolor{cgoLight}{HTML}{EEEEEE}
\definecolor{cgoExtralight}{HTML}{FFFFFF}

%
% Nuevo comando para tablas grandes sin cabecera.
\newcommand{\tablaSinCabeceraConBandas}[5]{%
  \begin{center}
    \tablefirsthead{
      \toprule
    }
    \tablehead{
      \multicolumn{#3}{l}{\small\sl continúa desde la página anterior}\\
      \hline
    }
    \tabletail{
      \hline
      \multicolumn{#3}{r}{\small\sl continúa en la página siguiente}\\
    }
    \tablelasttail{
      \hline
    }
    \bottomcaption{#1}
    \rowcolors[]{1}{cgoExtralight}{cgoLight}

  \begin{xtabular}{#2}
    #5
   \bottomrule
  \end{xtabular}
  \label{tabla:#4}
  \end{center}
}




\graphicspath{ {./img/} }

% Capítulos
\chapterstyle{bianchi}
\newcommand{\capitulo}[2]{
	\setcounter{chapter}{#1}
	\setcounter{section}{0}
	\setcounter{figure}{0}
	\setcounter{table}{0}
	\chapter*{#2}
	\addcontentsline{toc}{chapter}{#2}
	\markboth{#2}{#2}
}

% Apéndices
\renewcommand{\appendixname}{Apéndice}
\renewcommand*\cftappendixname{\appendixname}

\newcommand{\apendice}[1]{
	%\renewcommand{\thechapter}{A}
	\chapter{#1}
}

\renewcommand*\cftappendixname{\appendixname\ }

% Formato de portada
\makeatletter
\usepackage{xcolor}
\newcommand{\tutor}[1]{\def\@tutor{#1}}
\newcommand{\course}[1]{\def\@course{#1}}
\definecolor{cpardoBox}{HTML}{E6E6FF}
\def\maketitle{
  \null
  \thispagestyle{empty}
  % Cabecera ----------------
  \begin{minipage}{\textwidth}
    \begin{center}
      \includegraphics[width=\textwidth]{cabecera}\vspace{1cm}%
    \end{center}
  \end{minipage}
  \vfill
  % Título proyecto y escudo informática ----------------
  \colorbox{cpardoBox}{%
    \begin{minipage}{.8\textwidth}
      \vspace{.5cm}\Large
      \begin{center}
      \textbf{TFG del Grado en Ingeniería Informática}\vspace{.6cm}\\
      \textbf{\LARGE\@title{}}
      \end{center}
      \vspace{.2cm}
    \end{minipage}

  }%
  \hfill\begin{minipage}{.20\textwidth}
    \includegraphics[width=\textwidth]{escudoInfor}
  \end{minipage}
  \vfill
  % Datos de alumno, curso y tutores ------------------
  \begin{center}%
  {%
    \noindent\LARGE
    Presentado por \@author{}\\ 
    en Universidad de Burgos \\
    \@date{}\\
    Tutores: \nombreTutorJ\\
    \nombreTutorV\\
  }%
  \end{center}%
  \null
  \cleardoublepage
  }
\makeatother

\newcommand{\nombre}{Enrique Camarero Alonso}
\newcommand{\nombreTutores}{José Ignacio Santos Martín y Virginia Ahedo García}
\newcommand{\nombreTutorJ}{José Ignacio Santos Martín}
\newcommand{\nombreTutorV}{Virginia Ahedo García}

% Datos de portada
\title{GII 21.13 Sistema de recomendación de tecnologías para desarrolladores basado en la encuesta de Stack Overflow \\Documentación Técnica}
\author{\nombre}
\tutor{\nombreTutores}
\date{13 de junio de 2022}

\begin{document}

\maketitle

\cleardoublepage



\frontmatter


\clearpage

% Indices
\tableofcontents

\clearpage

\listoffigures

\clearpage

\listoftables

\clearpage

\mainmatter

\appendix

\apendice{Plan de Proyecto Software}

\section{Introducción}
En esta sección se va a comentar la evolución temporal del proyecto, la cual ha sido controlado mediante metodología ágil, y las viabilidades económicas y legales del proyecto.

\section{Planificación temporal}
El proyecto se divide en sprints y a la vez estos en tareas individuales. Todo esto es controlado desde un \href{https://github.com/eca1001/TFG}{repositorio creado específicamente para este proyecto en GitHub}. 

Se mostraran gráficas que muestren la duración de los sprints más relevantes. Esto se realizará a través de:

\begin{itemize}
    \item \textbf{Control chart}: esta gráfica muestra la duración de las tareas hasta su cierre y la media de estas durante todo el proyecto.
    \item \textbf{Cumulative flow}: esta gráfica muestra la evolución del estado de las tareas en cada momento, si están abiertas, se crean o se cierran.
\end{itemize}

\subsection{Sprint 0: Recogida de información (27/12/21 - 16/01/22)}
Este sprint consistía en familiarizarse con la encuesta de Stack Overflow y las herramientas que se iban a utilizar.

Las tareas que comprenden el sprint 0 son:
\begin{itemize}
    \item \textit{Networkx}: En esta tarea se iba a buscar información sobre Networkx, la creación de grafos y las posibilidades que nos ofrece al juntarlo con otras librerías como pandas.
    \item \textit{Generador HTML en Python}: Se empezaría a buscar como crear una página web desde Python para ejecutarlo en modo local.
    \item \textit{Sistema de evaluación del TFG}: Esta tarea se basa en la búsqueda de información acerca de todo lo que había que hacer y presentar (proyecto, memoria y anexos, en que se basaba la presentación, etc) al centro y al tribunal para poderme ir organizando y controlarlo.
    \item \textit{GitHub Desktop}: Se clonaría el proyecto en local con el fin de subir todos los cambios que se fueran realizando del proyecto al repositorio de GitHub.
\end{itemize}

\subsection{Sprint 1: Primeros pasos (17/1/22 - 29/1/22)}
En este sprint se va a iniciar el proyecto de Python y se va a tener el primer contacto con la encuesta a trabajar.

Las tareas correspondientes a este sprint son:
\begin{itemize}
    \item \textit{Importación de paquetes}: Se importaron todos aquellos paquetes con los que se iba a trabajar en las primeras semanas del proyecto.
    \item \textit{Descarga del CSV}: Se descargaría el archivo .csv que contiene los resultados de la encuesta de Stack Overflow.
    \item \textit{Primeros vistazos del CSV}: Se importó el archivo al código Python y se miraron los índices de las columnas con el fin de decidir con cual de ellas trabajar.
    \item \textit{Primeras redes}: Una vez se decidieron las columnas y con cual empezar (lenguajes con los que habían trabajado los encuestados), se crearon distintas redes sencillas, formadas por unos pocos encuestados, para ver su contenido . Una vez dado el visto bueno, se creó el primer grafo completo.
\end{itemize}

\subsection{Sprint 2: Primeros pasos (30/1/22 - 12/3/22)}
Se calculan las distintas métricas del grafo y sus nodos y se empieza a buscar una alternativa a Networkx para el mostrado del grafo, ya que este no permitía la interacción con el usuario.

Las tareas correspondientes a este sprint son:
\begin{itemize}
    \item \textit{Calculo de métricas}: Se calcularon métricas de los nodos como su influencia, su centralidad e importancia, cuales eran los vecinos más importantes de estos, y métricas de la red como la densidad del grafo.
    \item \textit{Búsqueda de mejores resultados}: Hasta ahora se había trabajado con el grafo al completo y los resultados no eran útiles, por lo que se empezó a buscar alternativas en busca de mejores resultados.
    \item \textit{Búsqueda de alternativas a Networkx para la generación de grafos}: Se empezó a investigar acerca de librerías que permitieran crear grafos interactivos y visualmente atractivos para el cliente.
\end{itemize}

La siguiente figura muestra el desarrollo de las issues en este sprint. \ref{fig:Sprint2}.

\imagen{Sprint2}{Gráfica control chart - Sprint 2}

\subsection{Sprint 3: Mejora de código 1/3/22 - 21/3/22}
Este sprint comenzaría antes de finalizar el anterior ya que se seguía buscando librerías alternativas a Networkx y no quería dejar parado el proyecto, ya que quería empezar a mejorar el código.

Este sprint esta formado por las siguientes tareas:
\begin{itemize}
    \item \textit{Comunidades}: Se creó un método que permite calcular las comunidades de nodos del grafo.
    \item \textit{Mejora de código}: Esta mejora se basó en crear un método que poda el grafo dependiendo de un umbral dictaminado por el usuario. De esta forma se cumpliría el requisito de que el cliente pudiera interactuar con el grafo adaptándolo a sus necesidades.
    \item \textit{Cursos básicos de HTML}: Empecé una serie de cursos para aprender el lenguaje y así poder utilizarlo.
    \item \textit{Estudio de Bokeh}: Me decidí por utilizar la librería Bokeh para sustituir a Networkx. De esta forma se cerraría el sprint anterior.
\end{itemize}

La siguiente imagen muestra la evolución de la apertura y cierre de las issues que componen los sprints hasta este momento. \ref{fig:Sprint0-3}.

\imagen{Sprint0-3}{Gráfica evolución - Sprint 0-3}



\subsection{Sprint 4: Creación página web 21/3/22 - 4/4/22}
Durante este tiempo se buscaron distintas maneras de crear una página web desde Python. A tavés de la librería webbrowser se encontró una solución. 
Este sprint está formado por el estudio, las pruebas y su posterior implementación de la librería en el proyecto, creando de esta forma una página web a la que podría añadir el grafo interactivo de Bokeh.

Una vez creada la página web en local, se desplegó en Visual Studio con el fin de conseguir una página editable gráficamente y fuera más profesional.

\subsection{Sprint 5: Despliegue en Heroku (21/3/22 - 6/4/22)}
Mientras se iba creando la aplicación web en Visual Studio, se estaba buscando la manera de desplegarla en un servidor web. El elegido fue Heroku.

Además, se añadieron cambios al código Python, añadiendo el método que crea el sistema de recomendación de tecnologías.

Las tareas de este sprint son:
\begin{itemize}
    \item \textit{Búsqueda de servidores web}: Se buscaron diferentes plataformas en la que desplegar la aplicación. Se buscó información sobre Plotly y Heroku, decantándome por la segunda. 
    
    \item \textit{Desplegar la aplicación en Heroku}: Inicialmente se desplegó el HTML que creaba Python ya que no se había finalizado la aplicación de Visual Studio para comprobar su funcionamiento.
    
    \item \textit{Opción de cambio de umbral en la aplicación}: Se ha añadido la opción al usuario de poder cambiar el umbral desde una pestaña en la aplicación.
    
    \item \textit{Creación del sistema de recomendación}: En el código Python, añadí un método que creará una tabla con los lenguajes y sus recomendaciones según el umbral seleccionado por el usuario.
    
    \item \textit{Guardar imágenes y grafos necesarios}: Creé en aquellos métodos que creaban imágenes para añadir al proyecto el comando que permite guardarlo en el sistema.
    
    \item \textit{Mejora de la presentación web}: Con la ayuda de CSS se mejoró el diseño de la página web.
\end{itemize}

La siguiente figura muestra el desarrollo de las issues en este sprint. \ref{fig:Sprint5}.

\imagen{Sprint5}{Gráfica control chart - Sprint 5}

\subsection{Sprint 6: Finalización página de lenguajes (15/4/22 - 29/4/22)}
Es este sprint se finalizó la página al completo de lenguajes.

Este sprint está formado por las siguientes tareas:
\begin{itemize}
    \item \textit{Rediseño del sistema de recomendación}: Tras una tutoría con los tutores se acordó la idea de cambiar el diseño, ya que se mostraba una tabla con todas las recomendaciones. En el nuevo diseño solo se muestran como máximo las 3 mejores.
    
    \item \textit{Creación de un filtro de búsqueda}: Al poder haber múltiples tecnologías que dificultaban la búsqueda de la buscada, creé un método que permite buscar la tecnología que se desea rápidamente.
    
    \item \textit{Creación tabla de tecnologías restantes}: Con la posibilidad de cambiar el umbral nos dimos cuenta que muchas tecnologías se perdían, por lo que decidí crear una tabla en la que aparecieran estas tecnologías junto un filtro de búsqueda.
    
    \item \textit{Unión de la página de lenguajes y el cambio de umbral}: Ya que el cambio de umbral se realizaba desde una pestaña aparte, decidí aumentar la usabilidad de la aplicación añadiendo el cambio de umbral en la misma página y acotándolo entre números enteros de 0 a 100.
    
    \item \textit{Cambio de nombres en la página}: El uso de nombres técnicos tanto en los gráficos como en el grafo podían dificultar a usuarios principiantes, por lo que estos nombres se cambiaron por unos menos complejos.
    
    \item \textit{Cambio de implementación en Heroku}: Con la página ya creada ya se podía cambiar la forma de desplegar la aplicación en Heroku a la versión creada en Visual Studio.
\end{itemize}

La siguiente figura muestra el desarrollo de las issues en este sprint. \ref{fig:Sprint6}.

\imagen{Sprint6}{Gráfica control chart - Sprint 6}


\subsection{Sprint 7: Crear resto de páginas de las tecnologías (24/4/22 - 3/5/22)}
Este sprint se basó en la creación del resto de las pestañas, las cuales estarían formadas por las distintas tecnologías.

Este sprint está compuesto por 7 tareas, siendo 6 de la creación de pestañas con cada una de las tareas:
\begin{itemize}
    \item \textit{Creación de pestañas}: Como se ha comentado, se crearon para cada tecnología una pestañas completa al igual que la de lenguajes creada en los anteriores sprints. Esto se ha realizado copiando el código que se tenía y adaptando los nombres del código a cada una de las tecnologías, creando en la template "base.html" cada una de las pestañas y en "main.py" el código para que redirecciones la pestaña a la página de esa tecnología.
    
    \item \textit{Juntar las pestañas de marcos de trabajo y otros marcos de trabajo en solo una pestaña}: Al implementar todas las tecnologías decidí que a pesar que en la encuesta los marcos de trabajo los separe en 2 apartados diferentes sería mejor si en el recomendador de lenguajes apareciesen todos juntos, ya que todos ellos forman parte de la misma tecnología. Se borró la página de otros marcos de trabajo y sus opciones de "base.html" y "main.py" y se creó una modificación en el código de marcos de trabajo para que relacionara cada uno de los marcos de trabajo tanto con los de su columna como con la otra.
\end{itemize}

La siguiente imagen muestra la evolución de la apertura y cierre de las issues que compone el sprint. \ref{fig:Sprint7}.

\imagen{Sprint7}{Gráfica evolución - Sprint 7}

\subsection{Sprint 8: Crear página de Inicio y finalización de la aplicación (4/5/22 - 11/5/22)}
En este último sprint se realizan 4 tareas muy diferenciadas que van a tener un fin común, la finalización de la aplicación:
\begin{itemize}
    \item \textit{Mejora del código de las tecnologías}: Se corrigió diferentes fallos que se encontraron en la creación de cada una de las páginas, como nombres mal escritos, llamadas a la red equivocada en las imágenes o grafos, variables, así como la eliminación de la variable umbral en ciertos métodos la cual no se utilizaba desde la versión de Jupyter Notebook.
    \item \textit{Finalización de la página de "Acerca de"}: Se finalizó la página que ayuda y explica fácilmente el proyecto y cómo se utiliza.
    \item \textit{No permitir no seleccionar ningún valor en el umbral}: Se producía un error si un usuario no rellenaba el formulario del umbral y ejecutaba la aplicación para ver los resultados. Esto se ha corregido obligando siempre al usuario a no dejar ese campo vacío si quiere ejecutar la página.
    \item \textit{Crear página de Inicio}: Se ha creado una página de inicio y que de esta forma sirva a modo de presentación del trabajo.
\end{itemize}

La siguiente figura muestra el desarrollo de las issues en este sprint. \ref{fig:Sprint8}.

\imagen{Sprint8}{Gráfica control chart - Sprint 8}

\section{Estudio de viabilidad}
Antes de empezar con el proyecto se va a estudiar la viabilidad tanto económica como legal que tendría el proyecto en el mercado, estimar sus costes y beneficios.

\subsection{Viabilidad económica}
En este apartado se realizará un estudio de los costes y beneficios del proyecto.

\subsubsection{Costes personales}
Se va a calcular el gasto que le supondría a la empresa de manera mensual y anual el llevar a cabo este proyecto.

El desarrollo del proyecto ha sido realizado solo por un alumno y dos profesores.

Se estima que se ha trabajado unas 600 horas durante 5 meses, es decir, unas 120 horas/mes aproximadamente. Por hora el alumno cobra 15€ brutos, lo que significa que cobraría 1800€/mes.

\tablaSmallSinColores{Costes salario alumno}{p{7cm} p{2cm} p{5cm}}{costes salario alumno}{
  \multicolumn{1}{p{2cm}}{\textbf{Concepto}} & \textbf{Coste (€)}\\
 }{
Salario mensual neto  & \multicolumn{1}{r}{1458}\\
Retención IRPF (12\%) & \multicolumn{1}{r}{216}\\
Seguridad social  (7\%) & \multicolumn{1}{r}{126}\\\hline
Salario mensual bruto  & \multicolumn{1}{r}{1800}\\\hline
\textbf{Salario total bruto(5 meses)}  & \multicolumn{1}{r}{9000}\\
}

En cuanto a los profesores digamos que tienen un salario de 30€/hora y trabajan 1 hora/semana, lo que supone 120€/mes brutos por profesor, a los cuales hay que sumar los impuestos.

\tablaSmallSinColores{Costes salario profesor}{p{7cm} p{2cm} p{5cm}}{costes salario profesor}{
  \multicolumn{1}{p{2cm}}{\textbf{Concepto}} & \textbf{Coste (€)}\\
 }{
Salario mensual neto  & \multicolumn{1}{r}{97.2}\\
Retención IRPF (12\%) & \multicolumn{1}{r}{14.4}\\
Seguridad social  (7\%) & \multicolumn{1}{r}{8.4}\\\hline
Salario mensual bruto  & \multicolumn{1}{r}{120}\\\hline
\textbf{Salario total bruto por profesor(5 meses)}  & \multicolumn{1}{r}{600}\\
}

Estos cálculos corresponden con el salario del alumno y los profesores, pero no con el coste que le supone a la empresa. A este valor hay que sumarle los impuestos que paga la empresa por él. Esto se puede consultar en \href{https://www.seg-social.es/wps/portal/wss/internet/Trabajadores/CotizacionRecaudacionTrabajadores/36537}{la página oficial de la seguridad social.} \cite{wiki:cotizacionSS}

Estos impuestos son los siguientes:
\begin{itemize}
    \item 23.6\% de contingencias
    \item 5.5\% de desempleo
    \item 0.2\% de FOGASA
    \item 0.6\% de formación profesional
\end{itemize}

\imagennormal{impuestosSS}{Régimen General de la Seguridad Social}

\newpage
Esto supone unos costes a la empresa de:
\tablaSmallSinColores{Costes del alumno a la empresa}{p{7cm} p{2cm} p{5cm}}{Costes del alumno a la empresa}{
  \multicolumn{1}{p{2cm}}{\textbf{Concepto}} & \textbf{Coste (€)}\\
 }{
Salario mensual bruto  & \multicolumn{1}{r}{1800}\\
Contingencias (23.6\%) & \multicolumn{1}{r}{424.8}\\
Desempleo  (5.5\%) & \multicolumn{1}{r}{99}\\
FOGASA  (0.2\%) & \multicolumn{1}{r}{3.6}\\
Formación  (0.6\%)  & \multicolumn{1}{r}{10.8}\\\hline
Coste total mensual  & \multicolumn{1}{r}{2338.2}\\\hline
\textbf{Coste total (5 meses)}  & \multicolumn{1}{r}{11691}\\
}

\tablaSmallSinColores{Costes de los profesores a la empresa}{p{7cm} p{2cm} p{5cm}}{Costes de los profesores a la empresa}{
  \multicolumn{1}{p{2cm}}{\textbf{Concepto}} & \textbf{Coste (€)}\\
 }{
Salario mensual bruto  & \multicolumn{1}{r}{120}\\
Contingencias (23.6\%) & \multicolumn{1}{r}{28.32}\\
Desempleo  (5.5\%) & \multicolumn{1}{r}{6.6}\\
FOGASA  (0.2\%) & \multicolumn{1}{r}{0.24}\\
Formación  (0.6\%)  & \multicolumn{1}{r}{0.72}\\\hline
Coste total mensual por profesor & \multicolumn{1}{r}{155.88}\\\hline
\textbf{Coste total por profesor(5 meses)}  & \multicolumn{1}{r}{779.4}\\
\textbf{Coste total ambos profesores(5 meses)}  & \multicolumn{1}{r}{1558.8}\\
}

\subsubsection{Costes hardware}
Se ha empleado un ordenador de sobremesa que se compró al principio del proyecto para la realización del mismo. Se supone que la amortización del mismo es de 5 años.

\tablaSmallSinColores{Costes hardware}{p{4cm} p{1.15cm} p{2.7cm}}{costes hardware}{
  \multicolumn{1}{p{4.5cm}}{\textbf{Concepto}} & \textbf{Coste} & \textbf{Amortización}\\
}{
Ordenador  & \multicolumn{1}{r}{800 €} & \multicolumn{1}{r}{67 €}\\
}

\subsubsection{Costes software}
Estos costes están compuestos por los costes asociados a las licencias software no gratuitas, en este caso la única licencia no gratuita es la del sistema operativo Windows 10 Home. Se supone que la amortización del mismo es de 5 años.

\tablaSmallSinColores{Costes software}{p{4cm} p{1.15cm} p{2.7cm}}{costes software}{
  \multicolumn{1}{p{4.5cm}}{\textbf{Concepto}} & \textbf{Coste} & \textbf{Amortización}\\
}{
Windows 10 Home  & \multicolumn{1}{r}{145 €} & \multicolumn{1}{r}{12 €}\\
}

El resto de las herramientas utilizadas para la realización el proyecto, como GitHub, Heroku o Visual Studio son de software libre y por tanto gratuitas. Se ha priorizado este tipo de herramientas sobre las de pago. 

\subsubsection{Costes totales}

\tablaSmallSinColores{Costes totales}{p{4cm} p{1.15cm} p{2.7cm}}{costes totales}{
  \multicolumn{1}{p{4.5cm}}{\textbf{Concepto}} & \textbf{Coste} & \textbf{Amortización}\\
}{
Alumno & \multicolumn{1}{r}{11691 €} & \multicolumn{1}{r}{-}\\
Profesores & \multicolumn{1}{r}{1558.8 €} & \multicolumn{1}{r}{-}\\
Ordenador  & \multicolumn{1}{r}{800 €} & \multicolumn{1}{r}{67 €}\\
Windows 10 Home  & \multicolumn{1}{r}{145 €} & \multicolumn{1}{r}{12 €}\\\hline
Total & \multicolumn{1}{r}{14194.8 €} & \multicolumn{1}{r}{79 €} \\
}

\subsubsection{Beneficios}
Este proyecto al ser un trabajo final de grado no está pensado para su comercialización, por lo que no se recibirían beneficios.

\subsection{Viabilidad legal}
Todo proyecto debe cumplir una serie de obligaciones legales. En este caso son las licencias software.

\tablaSmallSinColores{Viabilidad legal}{p{7cm} p{2cm} p{5cm}}{viabilidad legal}{
  \multicolumn{1}{p{2cm}}{\textbf{Software}} & \textbf{Licencia}\\
 }{
VirtualBox  & \multicolumn{1}{r}{Privativa / GPLv2}\\
Git  & \multicolumn{1}{r}{Privativa / GPLv2}\\
Heroku & \multicolumn{1}{r}{ISC}\\
}

Además, el proyecto ha utilizado la licencia MIT, que permite el uso comercial, modificación, distribución y uso privado.


\apendice{Especificación de Requisitos}

\section{Introducción}
La especificación de requisitos del software  ~\cite{wiki:ers} se basa en describir completamente el comportamiento del sistema a desarrollar, la interacción de los usuarios con el software a través de casos de uso y las necesidades tanto de clientes como usuarios. 

Su principal objetivo ~\cite{wiki:ersobjetivo} es funcionar como intermediario entre los usuarios, clientes y desarrolladores.

\section{Objetivos generales}
El objetivo principal del proyecto es la creación de un sistema de recomendación de tecnologías \cite{ekstrand2011} centrándose en los siguientes puntos:

\begin{itemize}
    \item Explotar los datos de la encuesta de de Stack Overflow mediante un sistema de recomendación basado en redes \cite{aggarwal2016}. Se va a recomendar distintas tecnologías entre las que se encuentras bases de datos, entornos de desarrollo, herramientas, lenguajes, marcos de trabajo y plataformas.
    \item Crear y diseñar una aplicación web que sea visualmente atractiva y fácil de comprender por todos los usuarios. Esta además debe ejecutarse en un servidor web para que se pueda acceder en cualquier momento y sin necesidad de descargarse nada.
    \item Cada una de las tecnologías deben tener su propia página web y deben tener un mismo patrón. Además, debe tener una página de inicio y otra de información en la que aparezca la información del alumno, tutores y una pequeña guía de uso.
    \item La aplicación debe ser interactiva, es decir, el usuario debe sentir que no está viendo una página web con información, sino que esa información se va cambiando a través de un valor introducido por el mismo usuario. Además, debe haber una red que muestre distintos datos de los nodos al colocarnos encima de cada uno de ellos.
\end{itemize}


\section{Catálogo de requisitos}
Se describirán los requisitos funcionales y los no funcionales del proyecto.

\subsection{Requisitos funcionales}
Son aquellos requisitos que el proyecto debe cumplir al finalizar su desarrollo.


  \begin{itemize}  
    \item \textbf{RF-1 Visualización de la información}: los datos deben ser claros y presentados de una manera atractiva para el usuario.
    \begin{itemize}
        \item \textbf{RF-1.1 Tabla de recomendaciones}: se añadirá una tabla con todos los nodos de la red junto con sus mejores recomendaciones por cada uno de éstos.
        \item \textbf{RF-1.2 Red interactiva}: se creará una red interactiva en la que, además de mostrar información de cada uno de los nodos al posicionarse encima de ellos, se deberán diferenciar las comunidades de nodos por colores.
        \item \textbf{RF-1.3 Histogramas y gráficas}: se añadirán histogramas y gráficas que muestren información acerca de la red de nodos.
    \end{itemize}
    
    \item \textbf{RF-2 Páginas independientes por tecnología}: cada tecnología deberá tener su propia pestaña con sus datos.
    
    \item \textbf{RF-3 Modificar umbral de poda}: la aplicación debe contener la posibilidad de variar el umbral de poda de la red.
    \begin{itemize}
        \item \textbf{RF-3.1 Modificación}: permitir al usuario cambiar el valor del umbral en cualquier momento.
        \item \textbf{RF-3.2 Verificación}: verificar que el campo no está vacío y el valor es un número entero entre 0 y 100.
    \end{itemize}
    
    \item \textbf{RF-4 Filtro de tecnologías}: cada una de las tablas de recomendaciones deberá tener un filtro de búsqueda.
    
    \item \textbf{RF-5 Página de inicio}: contendrá el título del proyecto, el escudo de la Universidad de Burgos, la Escuela Politécnica Superior y el del grado en ingeniería informática.
    
    \item \textbf{RF-6 Página de información}: la página contendrá tanto el nombre del autor como el de los tutores, así como un breve manual
    de uso de la aplicación.
    
\end{itemize}

\subsection{Requisitos no funcionales}
Son aquellos requisitos que no están relacionados con el uso del proyecto a través del usuario y los cuales se deben cumplir.

\begin{itemize}
    \item \textbf{RNF-1 Seguridad}: la aplicación no debe permitir modificaciones de código a ningún usuario que no sea el autor del proyecto.
    \item \textbf{RNF-2 Mantenibilidad}: el proyecto debe permitir modificaciones en el futuro de forma sencilla.
    \item \textbf{RNF-3 Portabilidad}: se debe permitir ejecutar la aplicación en cualquier plataforma y sistema operativo.
\end{itemize}

\section{Especificación de requisitos}

\subsection{Especificación de actores}
Existen 2 actores diferenciados:

\tablaSmallSinColores{Actores de la aplicación}{p{0.2\textwidth} | p{0.8\textwidth}}{Actores_aplicación}
{\textbf{Actor} & \textbf{Funcionalidad} \\}{
	Desarrollador & Se trata de la persona que ha desarrollado la aplicación. En este caso es \nombre. \\
	& El desarrollador estará ayudado por los tutores \nombreTutores para el desarrollo del proyecto. \\ \hline
	
	Cliente & Utilizan la aplicación. \\
}

\subsection{Diagramas de casos de uso}
En este apartado se mostrarán el diagrama de casos de uso. \ref{fig:CasoUso} \cite{wiki:casosDeUso}.

\imagen{CasoUso}{Diagrama de casos de uso}

\subsection{Especificación de casos de uso}

\tablaSmallSinColores{Caso de uso 1: Visualización de la información.}{p{3cm} p{.75cm} p{10cm}}{tablaCU1}{
	\multicolumn{3}{p{10.25cm}}{Caso de uso 1: Visualización de la información.} \\
}
{   
    Autor                                  & \multicolumn{2}{p{10.25cm}}{\nombre} \\\hline
    Versión                                & \multicolumn{2}{p{10.25cm}}{1.0} \\\hline
	Descripción                            & \multicolumn{2}{p{10.25cm}}{Visualizar el sistema de recomendación de las tecnologías.} \\\hline
	Precondiciones                         & \multicolumn{2}{p{10.25cm}}{Acceder a la aplicación} \\\hline
	Requisitos                         	   & \multicolumn{2}{p{10.25cm}}{RF-1.1, RF-1.2, RF-1.3} \\\hline
	\multirow{3}{3.5cm}{Secuencia normal}  & Paso & Acción \\\cline{2-3}
	& 1    & Acceder a la aplicación desde el navegador. \\\cline{2-3}
	& 2    & Navegar en las distintas pestañas que ofrece la aplicación. \\\hline
	Postcondiciones                        & \multicolumn{2}{p{10.25cm}}{Ninguna} \\\hline
	Excepciones                            & \multicolumn{2}{p{10.25cm}}{No tener acceso a internet.}\\\hline
	Frecuencia                             & Alta \\\hline
	Importancia                            & Alta \\\hline
	Urgencia                               & Alta \\
}

\tablaSmallSinColores{Caso de uso 2: Páginas independientes por tecnología.}{p{3cm} p{.75cm} p{10cm}}{tablaCU2}{
	\multicolumn{3}{p{10.25cm}}{Caso de uso 2: Páginas independientes por tecnología.} \\
}
{
    Autor                                  & \multicolumn{2}{p{10.25cm}}{\nombre} \\\hline
    Versión                                & \multicolumn{2}{p{10.25cm}}{1.0} \\\hline
	Descripción                            & \multicolumn{2}{p{10.25cm}}{Crear pestañas para cada una de las tecnologías.} \\\hline
	Precondiciones                         & \multicolumn{2}{p{10.25cm}}{Ninguna} \\\hline
	Requisitos                         	   & \multicolumn{2}{p{10.25cm}}{RF-2} \\\hline
	\multirow{3}{3.5cm}{Secuencia normal}  & Paso & Acción \\\cline{2-3}
	& 1    & Crear una pestaña completa para una de las tecnologías. \\\cline{2-3}
	& 2    & Replicar la pestaña creada para el resto de las tecnologías. \\\hline 
	Postcondiciones                        & \multicolumn{2}{p{10.25cm}}{Ninguna} \\\hline
	Excepciones                            & \multicolumn{2}{p{10.25cm}}{Ninguna}\\\hline
	Frecuencia                             & Alta \\\hline
	Importancia                            & Alta \\\hline
	Urgencia                               & Alta \\
}

\tablaSmallSinColores{Caso de uso 3: Modificar umbral de poda.}{p{3cm} p{.75cm} p{9cm}}{tablaCU3}{
	\multicolumn{3}{p{10.25cm}}{Caso de uso 3: Modificar umbral de poda.} \\
}
{
    Autor                                  & \multicolumn{2}{p{10.25cm}}{\nombre} \\\hline
    Versión                                & \multicolumn{2}{p{10.25cm}}{1.0} \\\hline
	Descripción                            & \multicolumn{2}{p{10.25cm}}{Modificar el valor del umbral con el que se poda la red.} \\\hline
	Precondiciones                         & \multicolumn{2}{p{10.25cm}}{Ninguna} \\\hline
	Requisitos                         	   & \multicolumn{2}{p{10.25cm}}{RF-3.1, RF-3.2} \\\hline
	\multirow{3}{3.5cm}{Secuencia normal}  & Paso & Acción \\\cline{2-3}
	& 1    & Cambiar el valor del umbral. \\\cline{2-3}
	& 2    & Pulsar la tecla \textit{Enter} o pinchar en \textit{Enviar}. \\\hline
	Postcondiciones                        & \multicolumn{2}{p{10.25cm}}{Ninguna} \\\hline
	Excepciones                            & \multicolumn{2}{p{10.25cm}}{Ninguna}\\\hline
	Frecuencia                             & Media \\\hline
	Importancia                            & Media \\\hline
	Urgencia                               & Baja \\
}

\tablaSmallSinColores{Caso de uso 4: Filtro de tecnologías.}{p{3cm} p{.75cm} p{9cm}}{tablaCU4}{
	\multicolumn{3}{p{10.25cm}}{Caso de uso 4: Filtro de tecnologías.} \\
}
{
    Autor                                  & \multicolumn{2}{p{10.25cm}}{\nombre} \\\hline
    Versión                                & \multicolumn{2}{p{10.25cm}}{1.0} \\\hline
	Descripción                            & \multicolumn{2}{p{10.25cm}}{Permite a los usuarios buscar rápidamente cada una de las tecnologías dentro de cada página.} \\\hline
	Precondiciones                         & \multicolumn{2}{p{10.25cm}}{Ninguna} \\\hline
	Requisitos                         	   & \multicolumn{2}{p{10.25cm}}{RF-4} \\\hline
	\multirow{3}{3.5cm}{Secuencia normal}  & Paso & Acción \\\cline{2-3}
	& 1    & Seleccionar la barra de escritura. \\\cline{2-3}
	& 2    & Escribir el nombre de la tecnología que se desea. \\\hline
	Postcondiciones                        & \multicolumn{2}{p{10.25cm}}{Ninguna} \\\hline
	Excepciones                            & \multicolumn{2}{p{10.25cm}}{Ninguna}\\\hline
	Frecuencia                             & Media \\\hline
	Importancia                            & Media \\\hline
	Urgencia                               & Baja \\
}

\tablaSmallSinColores{Caso de uso 5: Página de inicio.}{p{3cm} p{.75cm} p{9cm}}{tablaCU5}{
	\multicolumn{3}{p{10.25cm}}{Caso de uso 5: Página de inicio.} \\
}
{
    Autor                                  & \multicolumn{2}{p{10.25cm}}{\nombre} \\\hline
    Versión                                & \multicolumn{2}{p{10.25cm}}{1.0} \\\hline
	Descripción                            & \multicolumn{2}{p{10.25cm}}{Creación de una página de inicio de la aplicación.} \\\hline
	Precondiciones                         & \multicolumn{2}{p{10.25cm}}{Ninguna} \\\hline
	Requisitos                         	   & \multicolumn{2}{p{10.25cm}}{RF-5} \\\hline
	\multirow{3}{3.5cm}{Secuencia normal}  & Paso & Acción \\\cline{2-3}
	& 1    & Crear una página de inicio de la aplicación. \\\hline
	Postcondiciones                        & \multicolumn{2}{p{10.25cm}}{Ninguna} \\\hline
	Excepciones                            & \multicolumn{2}{p{10.25cm}}{Ninguna}\\\hline
	Frecuencia                             & Baja \\\hline
	Importancia                            & Baja \\\hline
	Urgencia                               & Baja \\
}

\tablaSmallSinColores{Caso de uso 6: Página de información.}{p{3cm} p{.75cm} p{9cm}}{tablaCU6}{
	\multicolumn{3}{p{10.25cm}}{Caso de uso 6: Página de información.} \\
}
{
    Autor                                  & \multicolumn{2}{p{10.25cm}}{\nombre} \\\hline
    Versión                                & \multicolumn{2}{p{10.25cm}}{1.0} \\\hline
	Descripción                            & \multicolumn{2}{p{10.25cm}}{Creación de una página de información de la aplicación.} \\\hline
	Precondiciones                         & \multicolumn{2}{p{10.25cm}}{Ninguna} \\\hline
	Requisitos                         	   & \multicolumn{2}{p{10.25cm}}{RF-6} \\\hline
	\multirow{3}{3.5cm}{Secuencia normal}  & Paso & Acción \\\cline{2-3}
	& 1    & Crear una página de información de la aplicación. \\\hline
	Postcondiciones                        & \multicolumn{2}{p{10.25cm}}{Ninguna} \\\hline
	Excepciones                            & \multicolumn{2}{p{10.25cm}}{Ninguna}\\\hline
	Frecuencia                             & Baja \\\hline
	Importancia                            & Baja \\\hline
	Urgencia                               & Baja \\
}

\apendice{Especificación de diseño}

\section{Introducción}
En este apéndice se van a detallar todos los aspectos relacionados con el diseño del proyecto.

\section{Diseño de datos}
Se ha utilizado la base de datos que Satck Overflow ofrece como resultado de la encuesta. La base de datos consta de 83440 filas y 48 columnas. Cada fila representa cada uno de las personas entrevistadas para la encuesta, y cada columna cada una de las preguntas.

De todas estas columnas se han elegido las siguientes:
\begin{itemize}
    \item Lenguajes con los que el usuario ha trabajado el último año.
    \item Bases de datos con las que el usuario ha trabajado el último año.
    \item Plataformas en la nube con las que el usuario ha trabajado el último año.
    \item Marcos de trabajo web con los que el usuario ha trabajado el último año.
    \item Herramientas con las que el usuario ha trabajado el último año.
    \item Entornos de desarrollo con los que el usuario ha trabajado el último año.
    \item Otros marcos de trabajo con los que el usuario ha trabajado el último año.
\end{itemize}


\section{Diseño procedimental}
En este apartado se va a describir el funcionamiento interno del proyecto mediante un diagrma de secuencias. \ref{fig:DiagramaSecuencias}.

\imagen{DiagramaSecuencias}{Diagrama de secuencias del funcionamiento interno}


\section{Diseño arquitectónico}
En este apartado se va describir los patrones que se han utilizado para el diseño del proyecto. \cite{wiki:patrones}

\subsection{Patrones estructurales}

\subsubsection{Fachada}
Este patrón pretende ofrecer al cliente una interfaz sencilla y disminuir el acoplamiento entre el subsistema y los clientes.

Se utiliza a la hora de presentar la aplicación y la comunicación de las páginas entre sí, ya que ofrece una interfaz sencilla que permite al usuario cambiar de pestañas y mantener el mismo umbral en todas ellas.

Su estructura es la siguiente: \ref{fig:Fachada}

\imagen{Fachada}{Patrón fachada}

\subsection{Patrones de comportamiento}
\subsubsection{Estado}
Este patrón permite a un objeto modificar su comportamiento cuando lo haga su estado interno.

En el proyecto se utiliza para crear nuevas redes en tiempo de ejecución una vez el usuario cambie el valor del umbral.

Su estructura es la siguiente: \ref{fig:Estado}

\imagen{Estado}{Patrón Estado}

\subsubsection{Observador}
Este patrón define una dependencia uno a varios para así en cuanto se modifique un estado todos sus dependientes sean notificados y actualizados automáticamente.

En el proyecto esto se utiliza para modificar las tablas, gráficos y redes que aparecen en cada una de las páginas. Una vez se modifica el umbral, se crea una nueva red y se notifica y actualiza en ese momento los objetos que aparecen en cada una de las páginas de las tecnologías. 

Su estructura es la siguiente: \ref{fig:Observador}

\imagen{Observador}{Patrón observador}

\subsubsection{Método plantilla}
Este patrón define el esqueleto de un algoritmo permitiendo a las subclases redefinir ciertos pasos del algoritmo sin cambiar su estructura general.

En el proyecto se utiliza para crear las páginas de las tecnologías. Las páginas heredan de una plantilla cambiando solamente en cada una de las tecnologías los datos a utilizar.

Su estructura es la siguiente: \ref{fig:Método plantilla}

\imagen{Método plantilla}{Patrón método plantilla}

\subsection{Patrones creacionales}
\subsubsection{Singleton}
Este patrón garantiza que una clase tenga tan solo una instancia, y que el acceso a la misma sea global. 

De esta forma al cliente tan solo le llega sólo un objeto (página web en este caso), y es este el que se modifica en cuanto se produce un cambio, no creando nuevas instancias.

Su estructura es la siguiente: \ref{fig:Singleton}

\imagen{Singleton}{Patrón singleton}
\apendice{Documentación técnica de programación}

\section{Introducción}
En este apéndice se va a describir la documentación técnica del proyecto, es decir, como está estructurado, como instalarlo y configurarlo correctamente, así como desplegarlo en Heroku y las pruebas que se han realizado.

\section{Estructura de directorios}
Esta estructura se puede encontrar en el código fuente de la aplicación, el cual se encuentra en GitHub.

\begin{itemize}
    \item \textbf{/}: es el directorio raíz. Aquí se encuentran aquellos ficheros necesarios para el despliegue de la aplicación en Heroku, la licencia, el archivo \textit{README} y el archivo \textit{main.py} desde el cual se ejecuta el proyecto.
    
    \item \textbf{/documentacion}: contiene la memoria y los anexos del proyecto.
    
    \item \textbf{/codigos}: en este directorio se guardan cada uno de los códigos de las tecnologías. Estos códigos se ejecutan desde el archivo \textit{main.py}.
    
    \item \textbf{/excel}: contiene dos carpetas:
    \begin{itemize}
        \item \textbf{vecinos}: son los archivos excel que se utilizaran para crear la tabla del sistema de recomendación.
        \item \textbf{faltantes}: son los archivos excel que se utilizaran para crear la tabla de los nodos eliminados tras la poda, los cuáles no aparecen ni en la red ni en los gráficos ni en el sistema de recomendación.
    \end{itemize}
    
    \item \textbf{/grafosTecnologias}: en este directorio se encuentran los grafos completos con los que se crearan las redes de cada una de las tecnologías.
    
    \item \textbf{/stack-overflow-developer-survey-2021}: contiene los archivos resultantes de la encuesta de Stack Overflow. En caso de que una tecnología no tenga grafo, se usarán estos resultados para crearlo.
    
    \item \textbf{/static}: contiene el archivo \textit{style.css} con el que se da estilo a la página web y las imágenes, gráficos, redes y tablas que se utilizan tanto para la página de inicio como en cada una de las páginas de las tecnologías.
    
    \item \textbf{/templates}: está compuesto por todas las páginas que se encargan del diseño gráfico de la página web.
    
\end{itemize}

\section{Manual del programador}
En este apartado se va a explicar qué programas han sido utilizados y cómo se han instalado para el desarrollo del proyecto.

El proyecto ha sido desarrollado en una máquina virtual Ubuntu de VirtualBox.

\subsection{Instalación de VirtualBox}
Para descargar VirtualBox tan solo hace falta ir a \href{https://www.virtualbox.org/wiki/Downloads}{su página web de descargas}. En esta página además de descargar el ejecutable para el sistema operativo en el que se va a ejecutar el programa (Windows e mi caso), será necesario descargar e instalar el paquete de extensión.\ref{fig:DescargaVBox}

\imagen{DescargaVBox}{Descarga de VirtualBox}

Una vez descargado, primero debemos instalar VirtualBox y después el paquete de expansión.

\subsection{Instalación máquina virtual Ubuntu}
Una vez descargado el sistema operativo en Windows, se procede a crear una nueva máquina virtual. Para ello se elige la opción dentro de VirtualBox. Se elige un nombre para esta nueva máquina, seleccionando la ubicación de la máquina, el tipo de sistema operativo y su versión, en este caso Linux Ubuntu versión 64 bits.

A continuación, hay que elegir la cantidad de memoria RAM que se reserva a esta máquina virtual, en este caso se ha elegido una de 8192 MB, sin embargo, lo recomendado para un correcto funcionamiento es 1024 MB. Después se elige la opción de crear un nuevo disco duro virtual para la nueva máquina. En este caso se elige crear un VDI que es una imagen de disco de VirtualBox para la máquina virtual. Se prefiere usar la opción de reservado dinámico de memoria ya que nos va a permitir usar la misma memoria, pero solo cuando se necesite y de esta forma la creación del disco es más rápida.

Posteriormente se escribe la ubicación dónde se almacenará el disco duro de la máquina 
virtual. Se selecciona el tamaño que tendrá el disco duro para el almacenamiento de los datos, 
en este caso 20 GB.

Este es el último paso, y como se puede ver en la siguiente imagen \ref{fig:VBoxConMaqVirt}, ya estaría lista la máquina virtual para ser ejecutada.

\imagen{VBoxConMaqVirt}{Máquina virtual creada}

\subsection{Configuración de la máquina virtual}
Antes de configurar la máquina virtual se debe instalar la imagen .iso del sistema operativo deseado, Ubuntu en este caso. Para ello nos dirigimos a \href{https://ubuntu.com/download/desktop}{su página oficial} y pinchamos en Download. \ref{fig:DescargaUbuntu}

\imagen{DescargaUbuntu}{Descarga de Ubuntu}

Una vez instalada la imagen .iso se ejecuta la máquina virtual anteriormente creada. A continuación, se selecciona la imagen de disco que previamente se había descargado de la página oficial de Ubuntu y se pincha en iniciar para que empiece la instalación.

\subsection{Instalación Visual Studio}
Para descargar Visual Studio nos dirigimos a \href{https://code.visualstudio.com/download}{su página oficial} y lo descargamos para el sistema operativo deseado. \ref{fig:DescargaVS}

\imagen{DescargaVS}{Descarga Visual Studio}


\section{Compilación, instalación y ejecución del proyecto}

\subsection{Descarga del proyecto}
El código fuente del proyecto se encuentra en el \href{https://github.com/eca1001/TFG}{repositorio de GitHub}.

Para descargarlo se debe pinchar en el botón verde con el combre \textit{Code} y seleccionar la opción de \textit{Download ZIP}. \ref{fig:DescargaRepositorio} Tras extraer el contenido del ZIP en la ruta deseada, vincularemos esta carpeta con el repositorio a través de GitHub Desktop, lo que va a permitir subir todos los cambios que se realicen en el proyecto a GitHub. 

\imagen{DescargaRepositorio}{Descarga del respositorio}

\subsection{Compilación y ejecución del proyecto}
Para compilar el proyecto se debe hacer mediante la terminal de Ubuntu. La primera vez que se vaya a compilar el proyecto se debe ejecutar los comandos para crear un entorno virtual:
\begin{itemize}
    \item \textit{sudo apt-get install python3-venv}
    \item \textit{python3 -m venv deploy}
\end{itemize}

Además, se deben tener instaladas las librerías que forman el archivo \textit{requirements.txt}, que son las que se utilizan en lo códigos de cada una de las tecnologías.

Una vez ya se ha instalado el entorno virtual para compilar el proyecto debemos acceder a la carpeta que contiene el proyecto y ejecutar los comandos: \ref{fig:CompilarProyecto}

\begin{itemize}
    \item \textit{. deploy/bin/activate}
    \item \textit{ cd deploy}
    \item \textit{ cd proyecto}
    \item \textit{python main.py}
\end{itemize}

\imagen{CompilarProyecto}{Comandos de compilación del proyecto}

Si no aparecen errores se puede ejecutar la aplicación en la dirección web \textit{http://127.0.0.1:5000}

\subsection{Despliegue en Heroku}
Para despleagar el proyecto en Heroku, lo primero que debemos hacer es registrarnos en \href{https://www.heroku.com/home}{su página web} y crear una aplicación desde el botón \textit{New}. Para crear la aplicación, como se ve en la siguiente imagen \ref{fig:CrearAppHeroku}, tan solo debemos darla un nombre y elegir una región entre Europa y Estados Unidos.

\imagen{CrearAppHeroku}{Crear aplicación en Heroku}

Una vez creada se entra en la aplicación y hay que dirigirse a la pestaña \textit{deploy}, donde veremos las opciones y los pasos a seguir para ejecutar el proyecto como se ve en la siguiente imahen \ref{fig:deployHeroku}.

\imagen{deployHeroku}{Desplegar la aplicación en Heroku}

En este momento la opción de desplegar a través de GitHub se encuentra desactivada, por lo que se ha realizado a través de Heroku CLI. Para ello primero se debe descargar, en caso de no tenerla, la herramienta Heroku CLI, la cuál podemos encontrar en \href{https://devcenter.heroku.com/articles/heroku-cli}{su página oficial}. Tras ello, basta con seguir los siguientes pasos:

\begin{itemize}
    \item \textit{heroku login} donde iniciamos sesión con nuestra cuenta.
    \item \textit{heroku git:remote -a nombreAplicación} donde nombreAplicación en mi caso es tfg-enriquecamareroalonso. Este comando sirve para conectarse con la aplicación en Heroku y así poder realizar commits con los cambios realizados.
    \item \textit{git add .}
    \item \textit{git commit -am "nombreDelCommit"}
    \item \textit{git push heroku main}
    \item 
\end{itemize}

Si no se ha producido ningún problema aparecerá un código parecido al siguiente en el cual se encuentra el enlace para ver la aplicación desde cualquier parte \ref{fig:ejecucionHeroku}.

\imagen{ejecucionHeroku}{Ejecución de comandos para despliegue en Heroku}


\section{Pruebas del sistema}
Para comprobar el correcto funcionamiento del proyecto se ha comprobado tanto por mi parte como la de los tutores, que todas las páginas funcionan correctamente, así como la opción del cambio del umbral.
\apendice{Documentación de usuario}

\section{Introducción}
En este apéndice se va a describir tanto los requisitos que deberán tener los usuarios para acceder a la aplicación como el manual de usuario de ésta.

\section{Requisitos de usuarios}
El único requisito que dispone la aplicación es la conexión a Internet, ya que se encuentra en un servidor web de Heroku y se accede mediante el siguiente enlace: \href{https://tfg-enriquecamareroalonso.herokuapp.com/}{tfg-enriquecamareroalonso.herokuapp.com/}

\section{Instalación}
No será necesario la instalación de ningún tipo para su uso.

\section{Manual del usuario}
En este apartado se detallará la utilización de la aplicación.

La aplicación tiene 3 partes diferenciadas: barra de navegación, contenido y pie de página. La primera y la última es común en todo el proyecto pero el contenido varia entre las distintas páginas.

\subsection{Barra de navegación}
La barra de navegación muestra las 8 pestañas por las que el usuario se puede mover en la página. Estas pestañas son: inicio, las 6 páginas de tecnologías (bases de datos, entornos de desarrollos, lenguajes, marcos de trabajo, herramientas, plataformas) y acerca de. \ref{fig:barraNavegacion}

Para desplazarse basta con pinchar sobre el nombre de la pestaña a la que se quiere ir.

\imagennormal{barraNavegacion}{Barra de navegación de la aplicación}

\subsubsection{Pie de página}
El pie de página muestra el nombre de la entidad en donde se ha realizado el proyecto (Universidad de Burgos), el tipo de proyecto (Trabajo Final de Grado) y el año (2022). \ref{fig:footer}

\imagennormal{footer}{Pie de página de la aplicación}

\subsection{Contenido}
Como antes se ha mencionado, el contenido depende de la pestaña en la que se encuentra el usuario.

\subsubsection{Página de Inicio}
El contenido de la página de inicio muestra el nombre del proyecto junto a los escudos de la Universidad de Burgos, la Escuela Politécnica Superior y el grado en ingeniería informática. \ref{fig:inicioPage}

\imagennormal{inicioPage}{Contenido de la página de inicio}

\subsubsection{Tecnologías}
Las páginas de tecnologías contienen el sistema de recomendación, además de una red interactiva de nodos y gráficas.

En el proyecto podemos ver que existen 6 páginas diferentes, una para cada tecnología con umbral a 30 por defecto:

\newpage
\textbf{Página de bases de datos}
\imagennormal{bbddPagina}{Página de bases de datos}

\newpage
\textbf{Entornos de desarrollo}
\imagennormal{entornosPagina}{Página de entornos de desarrollo}

\newpage
\textbf{Lenguajes}
\imagennormal{lenguajesPagina}{Página de lenguajes}

\newpage
\textbf{Marcos de trabajo}
\imagennormal{marcosPagina}{Página de marcos de trabajo}

\newpage
\textbf{Herramientas}
\imagennormal{herramientasPagina}{Página de herramientas}

\newpage
\textbf{Plataformas}
\imagennormal{plataformasPagina}{Página de plataformas}

\newpage
Si nos posicionamos encima de cada uno de los nodos en cualquiera de cada una de las pestañas, podemos observar como aparece un cuadro en el que aparece diferente información del nodo seleccionado.

\imagennormal{datosNodo30}{Datos nodo JavaScript con umbral a 30}

\newpage
Por otro lado, si introducimos información en el filtro del recomendador de tecnologías podemos observar como tan solo aparecen aquellas que contienen la cadena introducida en su nombre.

\imagennormal{recomendador}{Recomendador de lenguuajes}

Toda esta información que aparece en cada una de las páginas es variable, es decir, se puede modificar la información que aparece haciendo uso del umbral que se encuentra en la parte superior izquierda. Para modificarlo basta con introducir un número entero comprendido entre 0 y 100, inclusive. Se va a cambiar el umbral de 30 (por defecto) a 10.

\newpage
Como se puede observar en las siguientes imágenes, se han producido cambios respecto a la página de lenguajes previamente mostrada.

\imagennormal{lenguajesPagina10}{Página de lenguajes con umbral a 10}

Del mismo modo, la información de los nodos ha sido actualizada.
\imagennormal{datosNodo10}{Datos nodo JavaScript con umbral a 10}

\newpage
\subsubsection{Acerca de}
Muestra información como el nombre del autor, los tutores, el objetivo del proyecto y un pequeño manual de uso.
\imagennormal{AcercaDePage}{Página de información}



\bibliographystyle{plain}
\bibliography{bibliografiaAnexos}

\newpage\null\thispagestyle{empty}\newpage
\end{document}